In functional data analysis, discrete series are converted into functions
through a basis expansion, a step whose effect on subsequent analyses is rarely
examined in economic applications. This thesis compares B-spline, Fourier, and
wavelet (Haar, Daubechies~2, Symlet~4) bases for representing the annual gross
domestic product growth rates of 20 Latin American countries over 1960--2024,
extending the dataset of a previous B-spline study to include the COVID-19
contraction. The bases are compared through reconstruction error, the number of
nonzero coefficients at matched error, and functional principal component
analysis. The extended dataset reproduces the variance structure of the original
study (47.2\% against 47.1\% for the first three components). At matched error,
the adaptive B-spline uses fewer coefficients than the wavelets (19.5 knots
against approximately 38 of 128 coefficients), and the two families distribute
the error differently: B-splines follow large multi-year episodes but flatten
isolated shocks, whereas thresholded wavelets retain some isolated shocks but
shrink large episodes. The Fourier basis is the least suited, since the series
show no clear periodicity and its approximations oscillate over stable periods.
The wavelet energy decomposition identifies the time scales at which each
country's variation is concentrated. The number of principal components required
to explain 80\% of the variance ranges from three to eight depending on the basis
and smoothing level, showing that the choice of basis is a modelling decision
that conditions the results of functional principal component analysis.\\

\noindent \textbf{Keywords}:\\
Functional data analysis, basis expansion, wavelets, B-splines, functional
principal component analysis, economic growth, Latin America.
